% Page budget: 0.45 pages (conclusion) + 0.90 pages (discussion, from 07) | Owns: interpretation of every result,
%   limits of each claim, departures from the research plan, the answer to the research question.
% WRITING GUIDE
% Purpose: Explain the results that do not speak for themselves, state the limit of each claim, disclose the
%   departures from the research plan, and answer the research question per aspect.
% Start here: the frozen Results section (06) and the method sections it evaluates (03, 04).
% Reference check: introduce no new reference; reuse only verified sources already cited in the method sections.
% Cautions: no new results, runs or figures; every number from a float of 06, the appendix tables or
%   Thesis-writeup/Results/18-10-2026/CSV; untested explanations stated as candidates; no LPV vs LTI as a research
%   result; no predicted-bias comparison; no OBC network-overlap digit (R1 label mismatch).
% Section is finished when: every discussion point is tied to a float and states its limit, and the conclusion
%   answers each aspect and the research question in fewer words than the discussion.
% MUST ESTABLISH (Dirk 2026-10-18; handoff tasks/handoffs/2026-10-18-discussion-conclusion.md)
%   Form: one section "Discussion and Conclusion", discussion first, by aspect as numbered in the Introduction
%   (1 structure and training, 2 OBC interpretability, 3 generalisation), then limits and departures, then a short
%   conclusion in the same order. Results (06) is observational; explanations, limits and departures are here.
%   Interpretability = unique estimate of the identifiable combinations, not true-parameter recovery.
%   Style per paragraph: claim, evidence, limit, what would test it. Conclusion: one or two sentences per aspect,
%   one sentence of the central limit, one recommendation.
%   Checked 2026-10-18: the final runs build the encoder from data (D-271, logs thesis-results-4), so the old
%   "encoder from true parameters" point is no limitation; the K1-rate point of 07 is a diagnostic only (D-265);
%   the axis weighting point of 07 is resolved by W_r and NRMS_e selection (03, eq:aug_loss).
\section{Discussion and Conclusion}\label{sec:conclusion}
\ifdraft
\begin{itemize}
\item Discussion per aspect: explanation, limit and test of each result; then the limits and departures that bound all three (handoff 2026-10-18)
\item Conclusion: the answer per aspect and to the research question, the central limit, one recommendation
\end{itemize}
\fi
This section interprets the results of Section~\ref{sec:results} per aspect,
states the limits of each claim and the departures from the research plan,
and then answers the research question.

\subsection{Augmentation structure}
\ifdraft
\begin{itemize}
\item The gain with multisine is the absorber mode, carried by the network and the added states (Fig.~\ref{fig:res_frf}, Table~\ref{tab:res_obc}, Table~\ref{tab:app_structure})
\item Added states and PS are both needed; shown on validation with one seed (Table~\ref{tab:app_structure})
\item Without multisine the final model is slightly above the baseline; candidates only (Table~\ref{tab:res_structure}) (?)
\end{itemize}
\fi
The gain on the records with multisine comes from the absorber mode.
The final model has a resonance at 2.11 to 2.13\,$f_b$ over $Y$, against
2.12\,$f_b$ of the system (Fig.~\ref{fig:res_frf}).
Its damping ratio is 0.036 to 0.080 over $Y$, against 0.043, so the mode
frequency is learned more accurately than its damping.
With the learned block switched off, the $Y$ error of the trained model
returns to about 42\,\% (Table~\ref{tab:res_obc}), so the network and the
added states carry the mode, not the refitted parameters.
\\
Both the added states and their PS start are needed to learn the mode.
On the validation records with multisine, the $Y$ error is 6.4\,\% with PS,
78.1\,\% with plain training and 86.5\,\% without added states
(Table~\ref{tab:app_structure}).
Eight added states give the same error as two.
This rests on one seed per model, so a seed that lets plain training find
the mode is not excluded.
The missing dynamics hold one mode of known order, so the result shows that
two added states suffice for this mode, not how $n_{\bar x}$ should be
chosen when the order is unknown.
\\
Without multisine, the final model is 0.2 to 0.4 points above the baseline
at the true parameters (Table~\ref{tab:res_structure}).
At this level the baseline misses only the friction and the part of the
absorber that the motion excites.
A candidate explanation is that the learned correction, fitted mainly where
the multisine excites the absorber, adds small errors where the baseline is
already close to its limit.
The detuned start is not the cause: on the validation records, the model
trained from the true parameters is also above the baseline on $Y$
(Table~\ref{tab:app_structure}).
The candidate is not tested.
The error spectrum of both models on these records would show whether the
difference lies in the absorber band or below the crossover.
\todo{Missing: a check of the candidate. Candidate: compare the error spectra of the final model and the baseline on I3 to I6 with the existing evaluation code (no new run; basis: res\_ps\_spectrum uses the same computation on the validation multisine records).}

\subsection{Interpretability}
\ifdraft
\begin{itemize}
\item OBC halves the drift of U but does not recover the truth; expected from the truth overlap and \eqref{eq:obc_recovery} (Table~\ref{tab:res_obc})
\item The dampers move most; the friction projects onto the damping sensitivities by construction (D-204 argument, not measured)
\item OBC holds to first order at $\bar x=0$ for the one-step fit only; added states are not protected
\item Uniqueness only partly shown: the two starts still differ (Table~\ref{tab:app_obc_full}) (?: row requested)
\item Without added states the order reverses; candidate only (Table~\ref{tab:app_obc_na0})
\item Root cause analysis compares unique estimates only while the missing dynamics do not change
\end{itemize}
\fi
OBC reduces the parameter drift, but it does not recover the true
combinations.
From the detuned start, the combination error grows from 30.0\,\% to
47.7\,\% with OBC and to 83.1\,\% without (Table~\ref{tab:res_obc}).
From the correct start, OBC ends at 35.4\,\% and $\mathcal M_{\mathrm U}$ at
58.3\,\% (Table~\ref{tab:app_obc_full}).
This drift is expected from Section~\ref{sec:obc_ident}.
The truth overlap of 0.59 shows that $\Phi^\top\Xi\neq0$ on the tuples, so
the estimate of \eqref{eq:obc_recovery} is biased even with the projection.
OBC costs no accuracy, since the full models and their physical parts differ
by at most 0.9 points between the two arms.
\todo{Missing: the numbers of $\mathcal M_{\mathrm{OBC}}$ (R1, R2) are preliminary (update 19500 of 33400). Candidate: update 47.7, 35.4, 0.9 and the two-start differences after the final rerun (basis: handoff 2026-10-18, Sec.~3).}
\\
The damping combinations $c_{g1}$, $c_{g2}$ and $c_y$ move most, by 23 to
30 points with OBC and 43 to 46 points without, while the mass combinations
stay within 3.4\,\% of their true values.
A candidate cause is the rail friction of \eqref{eq:setup_fric}.
The sensitivity to $c_{g1}$ is proportional to $\dot X_1$, and the friction
force has the sign of $\dot X_1$.
Their inner product over a moving record is therefore nonzero, so the
friction lies partly in the damping directions and biases the dampers
through $\Phi^\top\Xi$.
This argument is ours and is not measured on the final data.
The truth overlap of $\Xi$ split into its friction and absorber parts would
test it.
\\
The guarantee of OBC is also narrower than the training it is applied to.
The projection holds to first order around $\beta_{\mathrm{lin}}$, for the
one-step fit on the reference slice $\bar x=0$
(Section~\ref{sec:obc_ident}).
With PS the added states are active, so the part of $f_{\mathrm{aug}}$ that
depends on $\bar x$ is not orthogonal to the parameter directions, and the
closed-loop criterion \eqref{eq:aug_loss} is not the one-step fit.
The two starts measure how far uniqueness holds in practice.
Their final estimates differ by 25.1\,\% with OBC and 36.4\,\% without, from
Table~\ref{tab:app_obc_full}.
OBC brings the two estimates closer, but they are not equal, so uniqueness is
only partly shown.
\todo{Missing: the two-start difference is computed by hand from Table~\ref{tab:app_obc_full}, not shown in a float. Candidate: the ``difference between starts'' row requested at the end of Section~\ref{sec:results}, computed by \texttt{res\_obc\_combinations.py} after the final rerun.}
\\
Without added states the order reverses: $\mathcal M_{\mathrm U}$ ends at
48.8\,\% and $\mathcal M_{\mathrm{OBC}}$ at 68.4\,\%
(Table~\ref{tab:app_obc_na0}).
A candidate explanation is that the network must then represent the absorber
statically, and a large part of it lies in the parameter directions (truth
overlap 0.59).
The projection removes that part from the network, so the parameters absorb
it.
This is not tested.
The result shows that OBC does not bring the estimate closer to the truth in
general.
Its guarantee is a unique split, and only the pair from two starts tests that.
\todo{Missing: a possible inconsistency in Table~\ref{tab:app_obc_na0}: its full-model $Y$ row (41.8, 42.7) equals the physical-part $Y$ row of Table~\ref{tab:res_obc}. Candidate: check \texttt{res\_obc\_na0} in the evaluation code before the paragraph above relies on that table.}
\\
Root cause analysis (Section~\ref{sec:intro}) compares unique estimates
over time.
By \eqref{eq:obc_recovery}, the estimate contains the bias
$(\Phi^\top\Phi)^{-1}\Phi^\top\Xi$, which depends on the missing dynamics.
A change of the estimate therefore indicates a changed physical parameter
only while the missing dynamics and the excitation stay the same.
This is our reasoning and is not tested.

\subsection{Generalisation}
\ifdraft
\begin{itemize}
\item The final model keeps its accuracy on each E record; each E record leaves the training range in one direction only (Fig.~\ref{fig:res_generalisation}, Fig.~\ref{fig:setup_coverage})
\item The black box does not fit the records without multisine, training records included; candidates, with the budget evidence (Table~\ref{tab:app_bb_search})
\item The comparison does not test the black box at its best; a black box of the plant inside the known loop would
\end{itemize}
\fi
The final model keeps its accuracy outside the training range.
Each extrapolation record is within 2.0 points of its interpolation record
(Fig.~\ref{fig:res_generalisation}).
On the multisine twins it stays below 10\,\% on $Y$, against 48 to 87\,\% for
the baseline.
Each extrapolation record leaves the training range in one direction only
($Y$, acceleration or velocity; Fig.~\ref{fig:setup_coverage}), so a
combined extrapolation is not tested.
The physical part extrapolates by construction through $M(Y)$, while the
network has only seen the training range.
\\
The black box does not fit the records without multisine, including the
training records, where its error is 72 to 106\,\%.
It therefore fails to fit these records rather than overfitting them.
It is below the final model only on $X_2$ of I2.
It models the whole closed loop from the reference and the injected force,
and must learn the rigid-body and controller dynamics from servo errors of
order $10^{-5}$~m.
Its selected configuration reached its best validation error at update 78000
of 80000 (Table~\ref{tab:app_bb_search}), so its training had not converged
within the budget.
Both are candidate causes and neither is tested.
A longer budget would test the second.
A black box of the plant inside the known loop would separate the first from
the comparison of the model classes.
The comparison therefore shows that this black box, at this budget, is less
accurate than the final model, not that every black box would be.

\subsection{Limits and departures from the research plan}
\ifdraft
\begin{itemize}
\item Simulation only, one absorber plus friction, one seed, the same known controller in data and model
\item Closed-loop training and evaluation can hide plant error behind the controller; no changed-controller test
\item Departures: LPV vs LTI aspect removed, controller transfer (E5) dropped, real Telica data out, settling dropped (D-241)
\end{itemize}
\fi
Three limits bound every result above.
The data are simulated, and the missing dynamics are one absorber mode and
rail friction, both chosen for this study.
Every model is trained once with one seed (Appendix~\ref{app:hyper}), so
differences of a few points between models are not resolved.
The controller in the model is the controller that generated the data.
\\
This last limit matters because training and evaluation are both closed
loop.
The controller corrects part of the output error, so a plant error at low
frequencies can stay hidden in the servo error.
An evaluation under a changed controller would expose it, and it is the use
that the Introduction motivates.
\\
The results depart from the research plan in four points.
The comparison of the LPV baseline with an LTI baseline is no longer an
aspect, because the LPV form follows from the equations of motion
(Section~\ref{sec:lfr}); it justifies the baseline in
Appendix~\ref{app:choices} (Fig.~\ref{fig:res_lpv_lti}).
The planned test with a changed controller is not in the results.
The planned evaluation on measured data from the industrial gantry is not
part of the results either.
Settling prediction is replaced by prediction over whole records, which the
metric of Section~\ref{sec:setup_metrics} evaluates.
\todo{Missing: a stated reason for dropping the changed-controller test (E5) and the measured data. Candidate: Dirk's decision with its basis (E5: not in the final evaluation, Introduction todo l.68; Telica: user decision 2026-09-24); the reason belongs here in one sentence each.}

\subsection{Conclusion}
\ifdraft
\begin{itemize}
\item Aspect 1: parallel augmentation with two added states and PS learns the missing mode from closed-loop data (Table~\ref{tab:res_structure})
\item Aspect 2: OBC halves the drift at no accuracy cost, gives a partly unique split, does not recover the truth
\item Aspect 3: the final model keeps its accuracy at extrapolation; better than this black box
\item Research question, central limit, one recommendation
\end{itemize}
\fi
A parallel augmentation with added states, initialised by PS and trained in
closed loop with the known controller, learns the missing absorber mode from
closed-loop data (Aspect~1).
With multisine it lowers the $Y$ error from 116.6\,\% to 6.2\,\%, and
without multisine it stays at the level of the baseline.
OBC halves the drift of the physical combinations at no cost in accuracy
(Aspect~2).
It makes the estimates from two starts closer but not equal, and it does
not recover the true values, because the missing dynamics are not orthogonal
to the parameter directions.
The augmented model keeps its accuracy at held-out positions, velocities and
accelerations, and is more accurate than the black box on every record
except one axis of one record (Aspect~3).
The answer to the research question is therefore a parallel augmentation
with added states, a PS start and the orthogonal construction: it predicts
the missing dynamics, and its parameters stay closer to a unique estimate
than without the construction.
The central limit is that uniqueness holds only to first order at
$\bar x=0$, and has been shown in simulation with one seed.
The next step is to evaluate the final model under a changed controller,
the use that motivates it.

% Source material kept until Dirk removes it (wip blocks of the earlier outline, 2026-10-11 to 2026-10-18).
\begin{wip}
\begin{itemize}
\item Answer whether LPV-LFR dynamic augmentation improves dual-gantry settling prediction over the physics baseline and black-box comparator
\item State whether the learned additional states contributed measurable missing dynamics under the final validated configuration
\item Conclude that state-level OBC prevents local tangent overlap by construction, while true parameters are recovered only when the physical discrepancy is orthogonal on the chosen excitation
\item Give the final quantitative trade between prediction error and physical combination error on held-out data
\item State whether the evidence covers relevant ASMPT motion, settling, controller transfer, or only a narrower simulation validation ladder
\item Close with the practical recommendation for when ASMPT should use the augmented model and when the baseline or a simpler residual model is safer
\item Figure plan: no conclusion figure; reserve the space for the final quantitative answer and limitations
\end{itemize}
\todo{Rewrite every conclusion bullet after freezing the final result table; no planned or diagnostic result may appear as achieved.}
\end{wip}
\begin{wip}
\begin{itemize}
    \item Why all states routing.
    \item why not use the encoder for the orthogonality each epoch instead of finite difference
    \item why not use the positions since we measure them --> train nn and encoder in combination? not full answer
    \item black box not fair comparison, no init (BLA) no controller
\end{itemize}

Thesis

Make a new folder for figures and code of the figures for the paper, and the final dataset that will be used.



* laten zien met parameters samenvoegen huidige methode niet correct ik moet iets van jacobian of iets laten zien
* ik heb nog steeds niet netjes uitgewerkt voor LPV-LFR met hoe ik het heb omgeschreven die n m ik kan niet zo goed zeggen hoe het heet
* verbeterde model met additional states laten zien als ik ze weg haal wat ze toevoegen = op nul zetten. ik weet niet of andere phase multisine validation en test genoeg laat zien hoe goed het werkt.
* Duidelijk verantwoorden controller. vooral met integrator problem van de plant en dus X en Y geen stijfheid K = 0 daar. kleine integratie error propegeert gwn door, dit zei maarten ook denk dat dit de error is.
* velocity central difference, why not get direct?
* Zijn Gramians nuttig om per state te laten zien?
* Multisine phase uitleggen kunnen we ook bronnen voor gebruiken. dat verschillende train validation en test verschillende multisine phase hebben.
* Met alleen simulatie is huidige simulatie niet genoeg, cubic spring toevoegen??
* benefit van mijn closed-loop training method is ook dat ik geen encoder / states er voor hoef te initializen. kessels wel referencen. figuur gebruiken.
* Jacobian additional states
* 2d / 3d grid wat jan zei van de data space of ik alles vang van het gebied van telica.
* Montecarlo simulatie zoals jan?
* For parameter recovery I had trajectories similar to the paper, currently not used for joint estimation?
* How well do we extrapolate outside multisine range? what does jan do?





Verdedigen

* waarom geen basis functions
* waarom niet zoals kessels: vooral andere literatuur kunnen zeggen waarom niet dat gebruikt, moet dat?
* wat is het voordeel van interpreteerbaarheid?--> controller maken, monitor health?
* waarom niet frf of mass lijn fitten voor de missende dynamica
* controller surpressses / hides issues such as orthogonality overlap. Byproduct of closed-loop.
* Hoe verdedig ik augmentation over residual frf fit??
* Hoeveelheid data kunnen verdedigen
* Huidige nf window and encoder window, should i run for different lengths?
* Orthogonality definition and that we orthogonalize additional states with respect to the baseline := 0. additional states propogate, not included in our definition.
* Kijken decisions.md
* do i need to justify controller can i say what i got from my supervisor, how to approach these kind of questions?
\end{wip}
